\documentclass[10pt]{article}
\usepackage[margin=0.55in,includeheadfoot]{geometry}
\usepackage[T1]{fontenc}
\usepackage[utf8]{inputenc}
\usepackage{lmodern}
\usepackage{multicol}
\usepackage{enumitem}
\usepackage{booktabs}
\usepackage{array}
\usepackage{tabularx}
\usepackage{amsmath,amssymb}
\usepackage{tikz}
\usepackage{graphicx}
\usepackage{eso-pic}
\usepackage{fancyhdr}
\usepackage{lastpage}
\usepackage{xcolor}
\usepackage[hidelinks]{hyperref}
\setlength{\columnsep}{0.28in}
\setlength{\parindent}{0pt}
\setlength{\parskip}{0.35em}
\raggedbottom
\pagestyle{fancy}
\fancyhf{}
\fancyhead[L]{\textsf{AP MICROECONOMICS}}
\fancyhead[C]{\textsf{Hard Unit Review Test Booklet}}
\fancyhead[R]{\textsf{Page \thepage\ of \pageref*{LastPage}}}
\fancyfoot[R]{\textsf{GO TO THE NEXT PAGE.}}
\renewcommand{\headrulewidth}{0.4pt}
\renewcommand{\footrulewidth}{0pt}
\newenvironment{choiceslist}
  {\begin{enumerate}[label=(\Alph*), leftmargin=1.6em, itemsep=0.12em, topsep=0.18em, parsep=0pt]}
  {\end{enumerate}}
\newenvironment{stimulus}
  {\begin{center}\begin{minipage}{0.96\columnwidth}\footnotesize\itshape\setlength{\tabcolsep}{3pt}\renewcommand{\arraystretch}{0.96}}
  {\end{minipage}\end{center}}
\newcommand{\unitheader}[2]{%
  \vspace{0.35em}
  {\large\bfseries Unit #1. #2\par}
  \vspace{0.2em}
  \hrule
  \vspace{0.45em}
}
\newcommand{\sourcebox}[1]{%
  \begingroup
  \small
  \setlength{\tabcolsep}{4pt}
  #1
  \endgroup
}
\newcommand{\CodexWatermark}{%
  \AddToShipoutPictureBG{%
    \AtPageCenter{%
      \makebox(0,0){%
        \rotatebox{45}{%
          \textcolor[gray]{0.92}{\fontsize{22}{22}\selectfont\bfseries Winston Huang-Generated with Codex}%
        }%
      }%
    }%
  }%
}
\begin{document}
\CodexWatermark
{\centering
{\LARGE\bfseries AP Microeconomics Intensive Unit 1--6 Review\par}
\vspace{0.5em}
{\large Hard Practice Test Booklet\par}
\vspace{1em}}
\textbf{Directions:} Choose the best answer for each question. The set is organized by unit, but the question style deliberately mixes calculation, graph interpretation, policy analysis, and multi-step reasoning in the style of difficult AP Microeconomics items. Some tables and figures were newly created for this booklet after synthesizing the recurring themes in the source PDFs from this directory.

\textbf{Structure:} 6 units, 20 questions per unit, 120 questions total.

\textbf{Reference base:} The booklet was synthesized from the TB/SG sets, FRQs, and progress checks found in this folder, especially the recurring topics on comparative advantage, elasticity, market intervention, perfect competition, imperfect competition, factor markets, externalities, public goods, and inequality.

\vspace{0.8em}
\begin{multicols*}{2}
\unitheader{1}{Basic Economic Concepts}
\begin{stimulus}
The following questions refer to Table 1.

\begin{center}
\begin{tabular}{lcc}
\toprule
Country & \shortstack{Hours per\\ Charger} & \shortstack{Hours per\\ Bottle} \\
\midrule
Alder & 2 & 6 \\
Briar & 3 & 4 \\
\bottomrule
\end{tabular}
\end{center}

Assume each country has 240 labor hours available and constant opportunity costs.
\end{stimulus}
\begin{enumerate}[series=apq, leftmargin=*, itemsep=0.85em, topsep=0.2em, parsep=0pt]
\item Which of the following statements is correct?
\begin{choiceslist}
\item Alder has an absolute advantage in both goods and a comparative advantage in solar chargers.
\item Briar has an absolute advantage in both goods and a comparative advantage in insulated bottles.
\item Alder has an absolute advantage in solar chargers, Briar has an absolute advantage in insulated bottles, and the comparative advantages are the same as the absolute advantages.
\item Briar has an absolute advantage in solar chargers, Alder has an absolute advantage in insulated bottles, and comparative advantage cannot be determined without demand data.
\item Neither country has an absolute advantage, but Alder has a comparative advantage in bottles.
\end{choiceslist}
\item If both countries specialize according to comparative advantage, the combined output will be
\begin{choiceslist}
\item 80 chargers and 40 bottles
\item 120 chargers and 60 bottles
\item 60 chargers and 120 bottles
\item 180 chargers and 40 bottles
\item 100 chargers and 80 bottles
\end{choiceslist}
\item Which of the following terms of trade for one solar charger would allow \emph{both} countries to gain from trade?
\begin{choiceslist}
\item $\tfrac{1}{5}$ of a bottle
\item $\tfrac{1}{3}$ of a bottle
\item $\tfrac{1}{2}$ of a bottle
\item $\tfrac{4}{5}$ of a bottle
\item 1 bottle
\end{choiceslist}
\item If the terms of trade are 1 charger for $\tfrac{1}{2}$ of a bottle and Alder exports 30 chargers, Alder will consume
\begin{choiceslist}
\item 60 chargers and 30 bottles, which is outside Alder's PPC
\item 90 chargers and 15 bottles, which is better than Alder could obtain on its own
\item 90 chargers and 15 bottles, which is worse than Alder could obtain on its own
\item 105 chargers and 10 bottles, which is exactly on Alder's PPC
\item 120 chargers and 0 bottles, which means Alder gains nothing from trade
\end{choiceslist}
\item Suppose Briar adopts a new technology that reduces the labor required per solar charger from 3 hours to 1.5 hours, while bottle technology is unchanged. Which of the following will occur?
\begin{choiceslist}
\item Briar gains both the absolute and comparative advantages in solar chargers.
\item Briar gains the absolute advantage in solar chargers, but Alder still has the comparative advantage in solar chargers.
\item Alder loses both the absolute and comparative advantages in solar chargers and bottles.
\item Comparative advantage disappears because Briar can now produce chargers more efficiently in absolute terms.
\item Briar gains the comparative advantage in bottles because its charger technology improved.
\end{choiceslist}
\end{enumerate}
\begin{stimulus}
The following questions refer to Table 2, which shows the daily production possibilities for drones and emergency medical kits in one economy.

\begin{center}
\begin{tabular}{ccc}
\toprule
Point & Drones & Medical Kits \\
\midrule
A & 0 & 300 \\
B & 10 & 294 \\
C & 20 & 282 \\
D & 30 & 258 \\
E & 40 & 216 \\
F & 50 & 150 \\
\bottomrule
\end{tabular}
\end{center}
\end{stimulus}
\begin{enumerate}[resume*=apq, leftmargin=*, itemsep=0.85em, topsep=0.2em, parsep=0pt]
\item The opportunity cost of increasing output from 20 drones to 30 drones is
\begin{choiceslist}
\item 6 medical kits
\item 12 medical kits
\item 18 medical kits
\item 24 medical kits
\item 42 medical kits
\end{choiceslist}
\item The data in Table 2 best illustrate
\begin{choiceslist}
\item constant opportunity costs because each additional drone costs the same number of kits
\item increasing opportunity costs because the number of kits forgone rises as more drones are produced
\item decreasing opportunity costs because drone specialization lowers the cost of each new drone
\item absolute advantage because the economy can produce both goods at the same time
\item comparative advantage because some points favor one good over the other
\end{choiceslist}
\item If the economy is currently producing 20 drones and 240 medical kits, which of the following is true?
\begin{choiceslist}
\item The economy is on its PPC and productively efficient.
\item The economy is inside its PPC, so it may be experiencing unemployment or inefficient resource use.
\item The economy is outside its PPC, so it must be trading internationally.
\item The economy is inside its PPC because opportunity costs are constant.
\item The economy is on its PPC because 240 kits lies between two listed output combinations.
\end{choiceslist}
\item If a new drone-specific technology raises the maximum number of drones from 50 to 70 while leaving the maximum number of medical kits at 300, the PPC will
\begin{choiceslist}
\item shift outward in a parallel way
\item pivot outward on the drone axis only
\item pivot inward on the medical-kit axis only
\item shift inward because resources become more specialized
\item become linear because technology changed only one sector
\end{choiceslist}
\item If the economy is currently at point E and policymakers require production of 10 additional drones with no technological change, the output of medical kits must fall by
\begin{choiceslist}
\item 24 kits
\item 30 kits
\item 42 kits
\item 50 kits
\item 66 kits
\end{choiceslist}
\end{enumerate}
\begin{stimulus}
The following questions refer to Lina's utility schedules. Lina buys only study guides and coffee drinks and spends all of her budget on those two goods.

\begin{center}
\begin{tabular}{cccc}
\toprule
\shortstack{Study\\ Guides} & MU & \shortstack{Coffee\\ Drinks} & MU \\
\midrule
1 & 30 & 1 & 18 \\
2 & 24 & 2 & 15 \\
3 & 18 & 3 & 12 \\
4 & 12 & 4 & 9 \\
5 & 6 & 5 & 6 \\
6 & 0 & 6 & 3 \\
\bottomrule
\end{tabular}
\end{center}

Each study guide costs \$6, each coffee drink costs \$3, and Lina's budget is \$24.
\end{stimulus}
\begin{enumerate}[resume*=apq, leftmargin=*, itemsep=0.85em, topsep=0.2em, parsep=0pt]
\item Which of the following bundles maximizes Lina's total utility?
\begin{choiceslist}
\item 1 study guide and 6 coffee drinks
\item 2 study guides and 4 coffee drinks
\item 3 study guides and 2 coffee drinks
\item 4 study guides and 0 coffee drinks
\item 0 study guides and 8 coffee drinks
\end{choiceslist}
\item Suppose Lina is currently buying 1 study guide and 6 coffee drinks. To increase total utility while staying within the same budget, she should
\begin{choiceslist}
\item buy more of both goods
\item buy fewer of both goods
\item buy more study guides and fewer coffee drinks
\item buy fewer study guides and more coffee drinks
\item keep the same bundle because the budget is fully exhausted
\end{choiceslist}
\item If Lina's budget and both prices double, Lina's utility-maximizing quantities will
\begin{choiceslist}
\item double
\item rise, but by less than double
\item fall because the goods become more expensive
\item remain unchanged
\item be impossible to determine without a demand curve
\end{choiceslist}
\item If the price of coffee drinks falls from \$3 to \$2 and all marginal utilities stay the same, which of the following is most likely?
\begin{choiceslist}
\item Lina will buy fewer coffee drinks because marginal utility diminishes.
\item Lina will buy more coffee drinks because their marginal utility per dollar rises.
\item Lina will buy the same quantities because the utility schedules did not change.
\item Lina will buy fewer study guides and fewer coffee drinks because the budget line rotates inward.
\item Lina will buy more study guides because coffee is now relatively cheaper.
\end{choiceslist}
\item The total utility Lina receives from her first 4 coffee drinks is
\begin{choiceslist}
\item 18 utils
\item 27 utils
\item 36 utils
\item 54 utils
\item 60 utils
\end{choiceslist}
\end{enumerate}
\begin{enumerate}[resume*=apq, leftmargin=*, itemsep=0.85em, topsep=0.2em, parsep=0pt]
\item Which of the following is most likely to shift an economy's future PPC outward by \emph{more} than an equal increase in current consumption would?
\begin{choiceslist}
\item Using more resources today to produce consumer goods rather than capital goods
\item Moving from a point on the PPC to a point inside the PPC
\item Devoting more current resources to capital formation and worker training
\item Imposing a binding price ceiling on capital goods
\item Specializing according to current comparative advantage only
\end{choiceslist}
\item Which of the following is a positive statement rather than a normative statement?
\begin{choiceslist}
\item The government should subsidize college because education is valuable.
\item A more equal income distribution is better for society than a less equal distribution.
\item A binding rent ceiling creates a shortage of apartments.
\item The minimum wage ought to be raised to a living wage.
\item Public parks are underprovided, so the government must build more of them.
\end{choiceslist}
\item A government wants to rely mostly on private markets but still correct obvious market failures and provide a social safety net. Which economic system best matches that description?
\begin{choiceslist}
\item A pure command economy
\item A mixed market economy
\item A traditional economy
\item A barter economy
\item A pure laissez-faire economy with no public sector
\end{choiceslist}
\item A firm paid \$500{,}000 for a nonrefundable feasibility study. The study now shows that producing the new product would add \$90{,}000 in marginal benefit and \$120{,}000 in marginal cost next year. The firm should
\begin{choiceslist}
\item produce the product because the study cost has already been incurred
\item produce the product because total spending on the project exceeds \$500{,}000
\item not produce the product because the additional cost exceeds the additional benefit
\item not produce the product only if average cost exceeds price
\item produce the product if and only if the study cost can be recovered later
\end{choiceslist}
\item 
                        A city is considering additional disaster-response drills. The marginal benefit and marginal cost of each additional drill are shown below.

                        \begin{center}
                        \begin{tabular}{ccc}
                        \toprule
                        Drill Number & Marginal Benefit & Marginal Cost \\
                        \midrule
                        1 & 18 & 5 \\
                        2 & 14 & 7 \\
                        3 & 10 & 9 \\
                        4 & 6 & 11 \\
                        5 & 2 & 13 \\
                        \bottomrule
                        \end{tabular}
                        \end{center}

                        To maximize net social benefit, the city should choose
                        
\begin{choiceslist}
\item 1 drill
\item 2 drills
\item 3 drills
\item 4 drills
\item 5 drills
\end{choiceslist}
\end{enumerate}
\unitheader{2}{Supply, Demand, and Market Intervention}
\begin{stimulus}
The following questions refer to the market for reusable water bottles.

\begin{center}
\begin{tabular}{ccc}
\toprule
Price & Quantity Demanded & Quantity Supplied \\
\midrule
\$4 & 96 & 24 \\
\$5 & 84 & 42 \\
\$6 & 72 & 60 \\
\$7 & 60 & 78 \\
\$8 & 48 & 96 \\
\bottomrule
\end{tabular}
\end{center}
\end{stimulus}
\begin{enumerate}[resume*=apq, leftmargin=*, itemsep=0.85em, topsep=0.2em, parsep=0pt]
\item In this market, the equilibrium price and quantity are
\begin{choiceslist}
\item \$5 and 42 units
\item \$5 and 84 units
\item \$6 and 60 units
\item \$6 and 72 units
\item \$7 and 78 units
\end{choiceslist}
\item If the government sets a binding price ceiling of \$5, the number of units exchanged in the market will be
\begin{choiceslist}
\item 24
\item 42
\item 60
\item 72
\item 84
\end{choiceslist}
\item If the government sets a binding price floor of \$7, the market will experience
\begin{choiceslist}
\item a shortage of 18 units
\item a shortage of 60 units
\item a surplus of 18 units
\item a surplus of 24 units
\item no inefficiency because the price floor is above equilibrium
\end{choiceslist}
\item Suppose reusable water bottles are a normal good and consumer incomes rise, while supply remains unchanged. Which of the following will occur?
\begin{choiceslist}
\item Equilibrium price will fall and equilibrium quantity will rise.
\item Equilibrium price will rise and equilibrium quantity will rise.
\item Equilibrium price will rise and equilibrium quantity will fall.
\item Equilibrium price will be unchanged and equilibrium quantity will rise.
\item Equilibrium price will be indeterminate and equilibrium quantity will rise.
\end{choiceslist}
\item Suppose production technology improves while, at the same time, consumer income rises and the good is inferior. Which of the following is definitely true?
\begin{choiceslist}
\item Equilibrium price will rise.
\item Equilibrium quantity will rise.
\item Equilibrium price will fall, but equilibrium quantity will be indeterminate.
\item Equilibrium quantity will fall, but equilibrium price will be indeterminate.
\item Both equilibrium price and equilibrium quantity will be indeterminate.
\end{choiceslist}
\end{enumerate}
\begin{stimulus}
The following questions refer to a ride-sharing service. At the current price, the price elasticity of demand for rides is 1.5, the cross-price elasticity of demand with bus rides is $+0.8$, and the income elasticity of demand is $+2.0$.
\end{stimulus}
\begin{enumerate}[resume*=apq, leftmargin=*, itemsep=0.85em, topsep=0.2em, parsep=0pt]
\item Based on these elasticity values, ride-sharing is best described as
\begin{choiceslist}
\item an inferior good and a complement to bus rides
\item a normal luxury good and a substitute for bus rides
\item a necessity and a complement to bus rides
\item a normal necessity good and an unrelated good with respect to bus rides
\item an inferior good and a substitute for bus rides
\end{choiceslist}
\item If the firm raises the price of rides by 10\%, total revenue from rides will most likely
\begin{choiceslist}
\item increase by 15\%
\item increase because demand is elastic
\item decrease because demand is elastic
\item remain unchanged because elasticity is constant
\item decrease only if supply is perfectly elastic
\end{choiceslist}
\item If bus fares fall by 5\%, the quantity demanded of ride-sharing services will most likely
\begin{choiceslist}
\item rise by 4\%
\item fall by 4\%
\item rise by 10\%
\item fall by 10\%
\item remain unchanged because cross elasticity applies only to complements
\end{choiceslist}
\item If household income falls by 3\%, the quantity demanded of ride-sharing services will most likely
\begin{choiceslist}
\item rise by 6\%
\item fall by 6\%
\item fall by 3\%
\item rise by 3\%
\item remain unchanged because the good is a luxury
\end{choiceslist}
\item If the demand for a product is less elastic than its supply, a per-unit excise tax on the product will be borne mostly by
\begin{choiceslist}
\item consumers, because the side of the market that is less responsive to price bears more of the tax
\item producers, because they remit the tax to the government
\item consumers and producers equally, because tax incidence is legally fixed
\item the government, because quantity traded falls
\item foreign producers, because imports absorb the tax
\end{choiceslist}
\end{enumerate}
\begin{stimulus}
The following questions refer to a market with demand curve $P = 30 - 0.5Q$ and supply curve $P = 6 + 0.25Q$, where $P$ is price and $Q$ is quantity.
\end{stimulus}
\begin{enumerate}[resume*=apq, leftmargin=*, itemsep=0.85em, topsep=0.2em, parsep=0pt]
\item In equilibrium, the price and quantity are
\begin{choiceslist}
\item \$12 and 24 units
\item \$14 and 32 units
\item \$16 and 28 units
\item \$18 and 24 units
\item \$20 and 20 units
\end{choiceslist}
\item At the free-market equilibrium, consumer surplus is
\begin{choiceslist}
\item \$128
\item \$192
\item \$224
\item \$256
\item \$288
\end{choiceslist}
\item If the government imposes a \$3 per-unit tax on sellers, the price paid by consumers will become
\begin{choiceslist}
\item \$13
\item \$14
\item \$15
\item \$16
\item \$17
\end{choiceslist}
\item With the \$3 per-unit tax in place, total tax revenue will be
\begin{choiceslist}
\item \$72
\item \$84
\item \$96
\item \$112
\item \$128
\end{choiceslist}
\item The deadweight loss created by the \$3 per-unit tax is
\begin{choiceslist}
\item \$4
\item \$6
\item \$8
\item \$12
\item \$16
\end{choiceslist}
\end{enumerate}
\begin{stimulus}
The following questions refer to a country that can trade coffee beans in the world market. Domestic demand is $Q_d = 120 - 2P$, domestic supply is $Q_s = 20 + 2P$, and the world price is \$15.
\end{stimulus}
\begin{enumerate}[resume*=apq, leftmargin=*, itemsep=0.85em, topsep=0.2em, parsep=0pt]
\item At the world price, the country will
\begin{choiceslist}
\item export 20 units
\item export 40 units
\item import 20 units
\item import 40 units
\item neither import nor export
\end{choiceslist}
\item Relative to autarky, opening trade at the world price of \$15 will
\begin{choiceslist}
\item benefit producers and harm consumers
\item benefit consumers and harm producers
\item benefit both consumers and producers
\item harm both consumers and producers
\item leave domestic welfare unchanged because the market clears either way
\end{choiceslist}
\item If the government imposes a \$5 tariff on imported coffee beans, the new domestic price will be
\begin{choiceslist}
\item \$15
\item \$18
\item \$20
\item \$25
\item \$30
\end{choiceslist}
\item With the \$5 tariff in place, tariff revenue will equal
\begin{choiceslist}
\item \$50
\item \$100
\item \$150
\item \$200
\item \$250
\end{choiceslist}
\item If the government replaced the tariff with an import quota of 20 units and all else remained the same, the domestic price would most likely
\begin{choiceslist}
\item remain at \$15 because quotas do not affect price
\item fall below \$15 because imports are restricted
\item rise to \$20 because the quota matches the post-tariff import quantity
\item rise to \$25 because the quota recreates autarky
\item be indeterminate because quotas affect only quantity
\end{choiceslist}
\end{enumerate}
\unitheader{3}{Production, Cost, and Perfect Competition}
\begin{stimulus}
The following questions refer to a firm's short-run production data. Labor is the only variable input, each worker is paid \$36 per day, and fixed cost is \$60.

\begin{center}
\begin{tabular}{cc}
\toprule
Workers & Total Product \\
\midrule
0 & 0 \\
1 & 18 \\
2 & 34 \\
3 & 48 \\
4 & 58 \\
5 & 64 \\
6 & 66 \\
\bottomrule
\end{tabular}
\end{center}
\end{stimulus}
\begin{enumerate}[resume*=apq, leftmargin=*, itemsep=0.85em, topsep=0.2em, parsep=0pt]
\item The marginal product of the fourth worker is
\begin{choiceslist}
\item 6 units
\item 10 units
\item 12 units
\item 14 units
\item 18 units
\end{choiceslist}
\item The law of diminishing marginal returns first appears when the firm hires the
\begin{choiceslist}
\item second worker
\item third worker
\item fourth worker
\item fifth worker
\item sixth worker
\end{choiceslist}
\item If the firm hires 5 workers, total variable cost will equal
\begin{choiceslist}
\item \$96
\item \$144
\item \$180
\item \$240
\item \$300
\end{choiceslist}
\item When output is 48 units, the firm's average fixed cost is
\begin{choiceslist}
\item \$0.75
\item \$1.00
\item \$1.25
\item \$1.50
\item \$2.00
\end{choiceslist}
\item The marginal cost of increasing output from 48 units to 58 units is
\begin{choiceslist}
\item \$2.40 per unit
\item \$3.60 per unit
\item \$6.00 per unit
\item \$10.00 per unit
\item \$36.00 per unit
\end{choiceslist}
\end{enumerate}
\begin{stimulus}
The following questions refer to a perfectly competitive firm with the total cost schedule below.

\begin{center}
\begin{tabular}{cc}
\toprule
Output & Total Cost \\
\midrule
0 & 30 \\
1 & 42 \\
2 & 50 \\
3 & 58 \\
4 & 70 \\
5 & 88 \\
6 & 114 \\
\bottomrule
\end{tabular}
\end{center}

Assume the market price is initially \$18 per unit.
\end{stimulus}
\begin{enumerate}[resume*=apq, leftmargin=*, itemsep=0.85em, topsep=0.2em, parsep=0pt]
\item The profit-maximizing output for this firm is
\begin{choiceslist}
\item 3 units
\item 4 units
\item 5 units
\item 6 units
\item the firm should shut down
\end{choiceslist}
\item At the profit-maximizing output, the firm's economic profit is
\begin{choiceslist}
\item \$0
\item \$2
\item \$10
\item \$18
\item \$20
\end{choiceslist}
\item The firm's shutdown price is closest to
\begin{choiceslist}
\item \$8.00
\item \$9.33
\item \$10.00
\item \$12.00
\item \$14.67
\end{choiceslist}
\item If the firm's fixed cost rises from \$30 to \$40, while market price remains \$18, the profit-maximizing output in the short run will
\begin{choiceslist}
\item fall to 4 units
\item fall to 3 units
\item remain at 5 units
\item rise to 6 units
\item fall to 0 because economic profit becomes negative
\end{choiceslist}
\item If the market price falls to \$10 per unit, the firm should
\begin{choiceslist}
\item shut down immediately because price is below average total cost
\item produce 2 units because price equals average variable cost there
\item produce 3 units because that is the highest output for which price still covers AVC and is near MC
\item produce 4 units because price is below marginal cost there
\item continue to produce 5 units to spread fixed cost over more output
\end{choiceslist}
\end{enumerate}
\begin{enumerate}[resume*=apq, leftmargin=*, itemsep=0.85em, topsep=0.2em, parsep=0pt]
\item A constant-cost perfectly competitive industry is in long-run equilibrium. If market demand increases, which of the following will occur in the long run?
\begin{choiceslist}
\item Price will remain above the original level and the number of firms will fall.
\item Price will return to the original level and total industry output will increase.
\item Price will remain above the original level and economic profit will remain positive.
\item Each firm's demand curve will shift rightward permanently.
\item The market will fail to return to zero economic profit.
\end{choiceslist}
\item A perfectly competitive firm is allocatively efficient when
\begin{choiceslist}
\item average total cost is minimized
\item price is greater than average total cost
\item price equals marginal cost
\item marginal revenue equals average variable cost
\item total revenue equals total cost
\end{choiceslist}
\item A perfectly competitive firm is productively efficient in the long run when
\begin{choiceslist}
\item it produces where marginal revenue exceeds marginal cost
\item it produces where average variable cost is maximized
\item it produces where average total cost is minimized
\item it produces where total revenue is maximized
\item it produces where price exceeds marginal cost
\end{choiceslist}
\item If a perfectly competitive market is currently experiencing economic losses for typical firms, the long-run adjustment will most likely involve
\begin{choiceslist}
\item entry of new firms and a lower market price
\item exit of firms and a higher market price
\item entry of new firms and a higher market price
\item exit of firms and a lower market price
\item no change, because perfectly competitive firms are price takers
\end{choiceslist}
\item Which of the following is true of the demand curve facing a single perfectly competitive firm?
\begin{choiceslist}
\item It slopes downward because the firm must lower price to sell more.
\item It is perfectly inelastic because the firm is small relative to the market.
\item It is perfectly elastic at the market price because the firm is a price taker.
\item It is the same as the market demand curve because the product is homogeneous.
\item It lies below the firm's marginal revenue curve.
\end{choiceslist}
\end{enumerate}
\begin{enumerate}[resume*=apq, leftmargin=*, itemsep=0.85em, topsep=0.2em, parsep=0pt]
\item A perfectly competitive firm is earning zero economic profit. If market demand for the product increases unexpectedly, the firm will most likely earn
\begin{choiceslist}
\item positive economic profit in the short run and zero economic profit again in the long run
\item positive economic profit in both the short run and the long run
\item zero economic profit in the short run and negative economic profit in the long run
\item negative economic profit in the short run because cost curves do not shift immediately
\item zero economic profit in both the short run and the long run because firms are price takers
\end{choiceslist}
\item If a competitive industry is an increasing-cost industry, a permanent increase in market demand will most likely cause the long-run equilibrium price to
\begin{choiceslist}
\item fall below the original price
\item remain equal to the original price
\item rise above the original price
\item be indeterminate
\item rise in the short run and then return below the original price
\end{choiceslist}
\item In the short run, a competitive firm should continue to produce even while earning losses if
\begin{choiceslist}
\item price is greater than or equal to average variable cost
\item price is greater than or equal to average total cost
\item marginal revenue is less than marginal cost
\item fixed cost is zero
\item economic profit is positive
\end{choiceslist}
\item If worker training raises the marginal product of labor for every worker in a perfectly competitive firm, then all else equal the firm's
\begin{choiceslist}
\item marginal cost curve will shift upward
\item marginal cost curve will shift downward
\item demand curve will shift upward
\item average fixed cost will rise
\item shutdown price will necessarily rise
\end{choiceslist}
\item A per-unit tax imposed on every firm in a constant-cost perfectly competitive industry will, in the long run, most likely
\begin{choiceslist}
\item leave the market price unchanged because firms are price takers
\item raise the market price by less than the tax and leave firms with losses forever
\item raise the market price enough to restore zero economic profit for firms
\item lower market price because industry supply becomes more elastic
\item cause firms to collude to offset the tax
\end{choiceslist}
\end{enumerate}
\unitheader{4}{Imperfect Competition}
\begin{stimulus}
The following questions refer to a single-price monopolist facing demand curve $P = 50 - Q$, marginal revenue curve $MR = 50 - 2Q$, constant marginal cost $MC = 10$, and constant average total cost $ATC = 18$.
\end{stimulus}
\begin{enumerate}[resume*=apq, leftmargin=*, itemsep=0.85em, topsep=0.2em, parsep=0pt]
\item The monopolist's profit-maximizing quantity is
\begin{choiceslist}
\item 10 units
\item 20 units
\item 25 units
\item 30 units
\item 40 units
\end{choiceslist}
\item The profit-maximizing price charged by the monopolist is
\begin{choiceslist}
\item \$18
\item \$20
\item \$25
\item \$30
\item \$40
\end{choiceslist}
\item The monopolist's economic profit is
\begin{choiceslist}
\item \$120
\item \$200
\item \$240
\item \$360
\item \$400
\end{choiceslist}
\item The allocatively efficient quantity in this market is
\begin{choiceslist}
\item 20 units
\item 25 units
\item 30 units
\item 40 units
\item 50 units
\end{choiceslist}
\item If technological improvement lowers marginal cost from 10 to 6, the monopolist will
\begin{choiceslist}
\item lower price and lower output
\item raise price and raise output
\item lower price and raise output
\item leave price unchanged and raise output
\item leave output unchanged and lower price
\end{choiceslist}
\end{enumerate}
\begin{stimulus}
Continue to assume the market demand is $P = 50 - Q$ and marginal cost is constant at \$10.
\end{stimulus}
\begin{enumerate}[resume*=apq, leftmargin=*, itemsep=0.85em, topsep=0.2em, parsep=0pt]
\item If the monopolist can perfectly price discriminate, the profit-maximizing quantity becomes
\begin{choiceslist}
\item 20 units
\item 25 units
\item 30 units
\item 40 units
\item 50 units
\end{choiceslist}
\item Compared with the single-price monopoly outcome, perfect price discrimination will
\begin{choiceslist}
\item reduce output and increase deadweight loss
\item increase output and eliminate deadweight loss
\item leave output unchanged and transfer producer surplus to consumers
\item lower output and eliminate consumer surplus
\item raise price for every unit sold
\end{choiceslist}
\item Under perfect price discrimination in this market, consumer surplus will be
\begin{choiceslist}
\item zero because the monopolist captures all willingness to pay above marginal cost
\item positive because quantity rises to the efficient level
\item equal to producer surplus by symmetry
\item larger than under single-price monopoly
\item negative because each consumer pays a different price
\end{choiceslist}
\item A movie theater charges students a lower ticket price than business travelers for the same film. This pricing policy is most likely based on the fact that
\begin{choiceslist}
\item students have more inelastic demand than business travelers
\item business travelers have more elastic demand than students
\item students generally have more elastic demand than business travelers
\item marginal cost of serving students is lower
\item students create larger positive externalities
\end{choiceslist}
\item Successful price discrimination generally requires all of the following EXCEPT
\begin{choiceslist}
\item some degree of market power
\item the ability to sort consumers into groups with different elasticities
\item the ability to prevent or limit resale
\item identical willingness to pay across all consumers
\item knowledge about demand differences across groups
\end{choiceslist}
\end{enumerate}
\begin{enumerate}[resume*=apq, leftmargin=*, itemsep=0.85em, topsep=0.2em, parsep=0pt]
\item Monopolistic competition is similar to perfect competition because both market structures have
\begin{choiceslist}
\item homogeneous products and significant barriers to entry
\item a downward-sloping demand curve for each firm
\item many firms and relatively free entry
\item mutual interdependence and strategic pricing
\item long-run prices above marginal cost
\end{choiceslist}
\item Monopolistic competition is similar to monopoly because firms in both market structures
\begin{choiceslist}
\item face a downward-sloping demand curve for their own product
\item have no close substitutes for their products
\item earn positive economic profit in the long run
\item produce where price equals marginal cost
\item face substantial legal barriers to entry
\end{choiceslist}
\item If a monopolistically competitive firm is earning positive economic profit in the short run, the most likely long-run adjustment is that
\begin{choiceslist}
\item its demand curve will shift right as brand loyalty rises
\item new firms will enter and its demand curve will shift left
\item existing firms will collude and reduce output
\item its marginal cost curve will shift upward until profit is zero
\item government will regulate price to average total cost
\end{choiceslist}
\item In long-run equilibrium, a monopolistically competitive firm typically operates
\begin{choiceslist}
\item at the minimum point of ATC and is productively efficient
\item where price equals marginal cost and is allocatively efficient
\item with excess capacity and output less than the quantity that minimizes ATC
\item where marginal revenue exceeds marginal cost
\item with a perfectly elastic demand curve
\end{choiceslist}
\item If advertising makes consumers view a firm's product as more differentiated, the firm's demand for its product will most likely become
\begin{choiceslist}
\item more elastic, reducing the firm's pricing power
\item less elastic, increasing the firm's pricing power
\item perfectly elastic, making the market more competitive
\item perfectly inelastic, allowing unlimited price increases
\item horizontal because more consumers know about the product
\end{choiceslist}
\end{enumerate}
\begin{stimulus}
The following questions refer to the payoff matrix below for two firms in an oligopoly, ColaCo and FizzCo. Each firm chooses either a High Price or a Low Price. Payoffs are in millions of dollars.

\begin{center}
\begin{tabular}{lcc}
\toprule
& FizzCo: High Price & FizzCo: Low Price \\
\midrule
ColaCo: High Price & (12,\ 12) & (4,\ 16) \\
ColaCo: Low Price & (16,\ 4) & (8,\ 8) \\
\bottomrule
\end{tabular}
\end{center}
\end{stimulus}
\begin{enumerate}[resume*=apq, leftmargin=*, itemsep=0.85em, topsep=0.2em, parsep=0pt]
\item ColaCo's dominant strategy is to
\begin{choiceslist}
\item choose High Price because 12 is greater than 8
\item choose High Price because collusion is more profitable
\item choose Low Price because it yields a higher payoff regardless of FizzCo's action
\item mix randomly between High Price and Low Price
\item there is no dominant strategy for ColaCo
\end{choiceslist}
\item The one-shot Nash equilibrium of this game is
\begin{choiceslist}
\item (High Price, High Price)
\item (High Price, Low Price)
\item (Low Price, High Price)
\item (Low Price, Low Price)
\item there is no Nash equilibrium
\end{choiceslist}
\item If the two firms could successfully collude and enforce the agreement, they would jointly choose
\begin{choiceslist}
\item (High Price, High Price) because joint profit is highest there
\item (Low Price, Low Price) because it is the Nash equilibrium
\item (Low Price, High Price) because total profit is 20
\item (High Price, Low Price) because total profit is 20
\item one firm would choose High Price and the other Low Price to maximize industry output
\end{choiceslist}
\item Why is collusion difficult to maintain in this oligopoly?
\begin{choiceslist}
\item Because each firm can raise joint profit by cutting price when the other keeps price high
\item Because the low-price strategy is illegal in oligopoly markets
\item Because firms in oligopoly face perfectly elastic demand curves
\item Because high price is a dominated strategy for both firms
\item Because consumers can force the collusive outcome through arbitrage
\end{choiceslist}
\item If government antitrust enforcement prevents explicit communication between the firms, the most likely outcome is that the market will
\begin{choiceslist}
\item move toward the low-price Nash equilibrium
\item guarantee the collusive outcome
\item become perfectly competitive overnight
\item eliminate strategic interdependence
\item force each firm to price at marginal cost
\end{choiceslist}
\end{enumerate}
\unitheader{5}{Factor Markets}
\begin{stimulus}
The following questions refer to a perfectly competitive firm that sells its output for \$12 per unit. The only variable input is labor, and the wage rate is initially \$60 per worker per day.

\begin{center}
\begin{tabular}{cc}
\toprule
Workers & Total Output \\
\midrule
0 & 0 \\
1 & 10 \\
2 & 18 \\
3 & 24 \\
4 & 28 \\
5 & 30 \\
\bottomrule
\end{tabular}
\end{center}
\end{stimulus}
\begin{enumerate}[resume*=apq, leftmargin=*, itemsep=0.85em, topsep=0.2em, parsep=0pt]
\item The marginal revenue product of the third worker is
\begin{choiceslist}
\item \$48
\item \$60
\item \$72
\item \$84
\item \$96
\end{choiceslist}
\item At the wage of \$60, the profit-maximizing number of workers hired is
\begin{choiceslist}
\item 1
\item 2
\item 3
\item 4
\item 5
\end{choiceslist}
\item If the market price of output rises from \$12 to \$15, the firm's demand for labor will
\begin{choiceslist}
\item shift left because labor is now more expensive
\item shift right because each worker's marginal revenue product rises
\item remain unchanged because the wage did not change
\item become perfectly elastic
\item shift left because diminishing returns set in sooner
\end{choiceslist}
\item If the wage falls from \$60 to \$45 while output price remains \$12, the firm will hire
\begin{choiceslist}
\item 2 workers
\item 3 workers
\item 4 workers
\item 5 workers
\item the same number of workers because labor demand depends only on product price
\end{choiceslist}
\item The firm's labor demand curve slopes downward primarily because of
\begin{choiceslist}
\item the law of diminishing marginal utility
\item the law of diminishing marginal product
\item the free-rider problem
\item price discrimination
\item economies of scale
\end{choiceslist}
\end{enumerate}
\begin{stimulus}
A bakery uses labor and ovens. At the current combination of inputs, the marginal product of labor is 24 loaves, the marginal product of ovens is 36 loaves, the wage is \$12, and the rental rate of an oven is \$18. The bakery sells bread in a perfectly competitive market for \$4 per loaf.
\end{stimulus}
\begin{enumerate}[resume*=apq, leftmargin=*, itemsep=0.85em, topsep=0.2em, parsep=0pt]
\item At the current combination of inputs, the bakery is
\begin{choiceslist}
\item using too much labor and too few ovens
\item using too few labor hours and too many ovens
\item using a cost-minimizing combination of labor and ovens
\item using an allocatively inefficient combination of inputs because output price exceeds wage
\item unable to determine whether the mix is cost-minimizing without total cost data
\end{choiceslist}
\item If the wage rises to \$16 while all else remains the same, to restore least-cost production the bakery should
\begin{choiceslist}
\item use more labor and fewer ovens
\item use fewer labor hours and more ovens
\item use more of both inputs
\item use less of both inputs
\item make no change because only relative productivity matters
\end{choiceslist}
\item If the marginal product of ovens rises from 36 to 54 while input prices are unchanged, the bakery should
\begin{choiceslist}
\item use more labor and fewer ovens
\item use fewer labor hours and more ovens
\item use less of both inputs
\item make no change because the firm is already cost minimizing
\item raise output price to restore equality between the two ratios
\end{choiceslist}
\item If the price of bread falls from \$4 to \$3, the bakery's least-cost combination of labor and ovens will
\begin{choiceslist}
\item shift toward labor because labor is cheaper than ovens
\item shift toward ovens because marginal revenue product falls
\item remain unchanged because least-cost input choice depends on input prices and physical productivity, not output price
\item be impossible to determine because bread demand also changes
\item require less of both inputs at every output level and therefore change the input ratio
\end{choiceslist}
\item At the current levels of input use, the marginal revenue product of labor is
\begin{choiceslist}
\item \$12, so the bakery is using labor efficiently
\item \$24, so the bakery should reduce labor
\item \$48, so the bakery should hire more labor
\item \$72, so the bakery should hire more labor
\item \$96, so the bakery should reduce labor
\end{choiceslist}
\end{enumerate}
\begin{enumerate}[resume*=apq, leftmargin=*, itemsep=0.85em, topsep=0.2em, parsep=0pt]
\item If the demand for smartphones increases, the demand for assembly-line labor used to produce smartphones will most likely
\begin{choiceslist}
\item shift left because labor is a normal good
\item shift right because labor demand is derived from product demand
\item become perfectly elastic
\item remain unchanged unless wages change first
\item shift left because the price of the final product rises
\end{choiceslist}
\item If workers receive training that makes them more productive, the demand curve for those workers will most likely
\begin{choiceslist}
\item shift left because their wage will rise
\item shift right because their marginal revenue product increases
\item become perfectly inelastic
\item shift left because the supply of skilled labor rises
\item remain unchanged because training affects only supply
\end{choiceslist}
\item In a perfectly competitive labor market, an effective minimum wage above equilibrium will most likely create
\begin{choiceslist}
\item no unemployment because firms must hire at the legal wage
\item excess demand for labor
\item excess supply of labor
\item a rightward shift of labor demand
\item a leftward shift of labor supply
\end{choiceslist}
\item If a binding minimum wage increases the total earnings of labor, that outcome is most likely because labor demand is
\begin{choiceslist}
\item relatively elastic
\item relatively inelastic
\item perfectly elastic
\item unit elastic
\item perfectly inelastic
\end{choiceslist}
\item If labor and capital are substitutes in production, an increase in the rental price of capital will most likely
\begin{choiceslist}
\item decrease the demand for labor
\item increase the demand for labor
\item decrease the supply of labor
\item increase the supply of labor
\item have no effect on labor demand
\end{choiceslist}
\end{enumerate}
\begin{stimulus}
The following questions refer to a monopsonistic labor market with labor supply curve $w = 10 + 2L$, marginal factor cost $MFC = 10 + 4L$, and marginal revenue product $MRP = 70 - 2L$, where $w$ is the wage and $L$ is the number of workers.
\end{stimulus}
\begin{enumerate}[resume*=apq, leftmargin=*, itemsep=0.85em, topsep=0.2em, parsep=0pt]
\item The monopsonist will hire
\begin{choiceslist}
\item 10 workers at a wage of \$30
\item 10 workers at a wage of \$50
\item 15 workers at a wage of \$40
\item 15 workers at a wage of \$30
\item 20 workers at a wage of \$50
\end{choiceslist}
\item In a competitive labor market with the same supply and demand conditions, employment and the wage would be
\begin{choiceslist}
\item 10 workers and \$30
\item 10 workers and \$40
\item 15 workers and \$40
\item 15 workers and \$50
\item 20 workers and \$40
\end{choiceslist}
\item Compared with the competitive outcome, the monopsony results in
\begin{choiceslist}
\item higher wages and higher employment
\item higher wages and lower employment
\item lower wages and lower employment
\item lower wages and higher employment
\item the same wage but lower employment
\end{choiceslist}
\item If the government sets a minimum wage of \$40 in this monopsony, employment will most likely
\begin{choiceslist}
\item fall to 10 workers
\item rise to 15 workers
\item rise above 15 workers
\item remain unchanged at 10 workers
\item be impossible to determine without the firm's total cost curve
\end{choiceslist}
\item The reason the monopsonist's marginal factor cost curve lies above the labor supply curve is that
\begin{choiceslist}
\item the firm pays each additional worker a higher wage but leaves existing workers' wages unchanged
\item the firm must raise the wage on all workers when it hires an additional worker
\item the labor supply curve is perfectly elastic
\item marginal revenue product rises with employment
\item fixed costs are increasing
\end{choiceslist}
\end{enumerate}
\unitheader{6}{Market Failure and the Role of Government}
\begin{stimulus}
The following questions refer to a competitive market in which demand is $P = 50 - Q$, the marginal private cost is $MPC = 10 + Q$, and the marginal social cost is $MSC = 20 + Q$.
\end{stimulus}
\begin{enumerate}[resume*=apq, leftmargin=*, itemsep=0.85em, topsep=0.2em, parsep=0pt]
\item The quantity produced in the unregulated market equilibrium will be
\begin{choiceslist}
\item 10 units
\item 15 units
\item 20 units
\item 25 units
\item 30 units
\end{choiceslist}
\item The socially optimal quantity is
\begin{choiceslist}
\item 10 units
\item 15 units
\item 20 units
\item 25 units
\item 30 units
\end{choiceslist}
\item A Pigouvian tax that would lead the market to the socially optimal quantity should be set at
\begin{choiceslist}
\item \$5 per unit
\item \$10 per unit
\item \$15 per unit
\item \$20 per unit
\item \$25 per unit
\end{choiceslist}
\item The deadweight loss from the unregulated market is
\begin{choiceslist}
\item \$12.50
\item \$25.00
\item \$37.50
\item \$50.00
\item \$75.00
\end{choiceslist}
\item The market described above is inefficient because it
\begin{choiceslist}
\item underproduces relative to the socially optimal quantity
\item overproduces because firms ignore external costs
\item underprices output because demand is too elastic
\item fails only because the market is perfectly competitive
\item overproduces because marginal social benefit exceeds demand
\end{choiceslist}
\end{enumerate}
\begin{stimulus}
The following questions refer to a flood barrier, which is a pure public good. Three residents have marginal benefit schedules given by
$MB_A = 30 - Q$, $MB_B = 24 - Q$, and $MB_C = 18 - Q$.
The marginal cost of each unit of the barrier is constant at \$36.
\end{stimulus}
\begin{enumerate}[resume*=apq, leftmargin=*, itemsep=0.85em, topsep=0.2em, parsep=0pt]
\item The marginal social benefit curve for the flood barrier is
\begin{choiceslist}
\item $MSB = 30 - Q$
\item $MSB = 72 - Q$
\item $MSB = 72 - 3Q$
\item $MSB = 108 - 3Q$
\item $MSB = 36 - 3Q$
\end{choiceslist}
\item The efficient quantity of the flood barrier is
\begin{choiceslist}
\item 8 units
\item 10 units
\item 12 units
\item 18 units
\item 36 units
\end{choiceslist}
\item If the flood barrier were left entirely to private voluntary purchase, the most likely outcome would be
\begin{choiceslist}
\item overproduction because people double count the benefits
\item underproduction because of the free-rider problem
\item efficient production because price reveals all benefits
\item production at the monopolistic quantity
\item production at the allocatively efficient quantity because the good is nonrival
\end{choiceslist}
\item A pure public good is characterized by
\begin{choiceslist}
\item rivalry and excludability
\item rivalry and nonexcludability
\item nonrivalry and excludability
\item nonrivalry and nonexcludability
\item positive externalities and rivalry
\end{choiceslist}
\item Which of the following is the best example of a common resource rather than a public good?
\begin{choiceslist}
\item National defense
\item A lighthouse signal
\item A fish population in an open-access lake
\item A citywide tornado siren
\item Broadcast radio
\end{choiceslist}
\end{enumerate}
\begin{enumerate}[resume*=apq, leftmargin=*, itemsep=0.85em, topsep=0.2em, parsep=0pt]
\item If a natural monopoly is regulated so that price equals marginal cost, the firm will often need a subsidy because
\begin{choiceslist}
\item marginal cost pricing always gives the firm economic profit
\item at the allocatively efficient output, price can lie below average total cost
\item natural monopolies have perfectly elastic demand curves
\item average fixed cost is zero in the long run
\item marginal revenue exceeds price for monopolies
\end{choiceslist}
\item If a natural monopoly is instead regulated at the price where demand intersects average total cost, the result is that the firm
\begin{choiceslist}
\item earns losses and produces the allocatively efficient quantity
\item earns zero economic profit but produces less than the allocatively efficient quantity
\item earns monopoly profit and produces more than the efficient quantity
\item must shut down because marginal cost exceeds price
\item becomes perfectly competitive
\end{choiceslist}
\item Antitrust policy is primarily intended to
\begin{choiceslist}
\item guarantee equal incomes across households
\item eliminate all large firms regardless of efficiency
\item limit anticompetitive practices such as collusion and price fixing
\item force all firms to produce where average total cost is minimized
\item replace markets with government planning
\end{choiceslist}
\item A cap-and-trade program for pollution can reduce total abatement cost relative to a uniform command-and-control rule because it
\begin{choiceslist}
\item lets firms with lower abatement costs reduce more pollution and sell permits
\item requires every firm to cut pollution by the same percentage
\item eliminates the need to define property rights
\item guarantees zero pollution in equilibrium
\item works only when demand is perfectly inelastic
\end{choiceslist}
\item In a monopsonistic labor market, a minimum wage set between the monopsony wage and the competitive wage can
\begin{choiceslist}
\item reduce both wages and employment
\item increase wages but always decrease employment
\item increase both wages and employment
\item have no effect because the monopsonist is already efficient
\item eliminate marginal factor cost
\end{choiceslist}
\end{enumerate}
\begin{stimulus}
The following table shows cumulative shares of national income for two countries.

\begin{center}
\begin{tabular}{lcc}
\toprule
\shortstack{Cumulative\\ Households} & \shortstack{Country A\\ Income Share} & \shortstack{Country B\\ Income Share} \\
\midrule
Bottom 20\% & 4\% & 8\% \\
Bottom 40\% & 13\% & 23\% \\
Bottom 60\% & 29\% & 43\% \\
Bottom 80\% & 53\% & 68\% \\
Bottom 100\% & 100\% & 100\% \\
\bottomrule
\end{tabular}
\end{center}
\end{stimulus}
\begin{enumerate}[resume*=apq, leftmargin=*, itemsep=0.85em, topsep=0.2em, parsep=0pt]
\item Based on the table, income is more equally distributed in
\begin{choiceslist}
\item Country A, because its top quintile receives a smaller share of income
\item Country A, because its bottom quintile receives a smaller share of income
\item Country B, because its Lorenz curve lies closer to the 45-degree line of equality
\item Country B, because its top quintile receives a larger share of income
\item neither country, because both have five quintiles
\end{choiceslist}
\item Which of the following is most likely true about the Gini coefficients of the two countries?
\begin{choiceslist}
\item Country A has the smaller Gini coefficient.
\item Country B has the smaller Gini coefficient.
\item The two countries must have the same Gini coefficient because total income is 100\% in both.
\item Country B's Gini coefficient must exceed 1.
\item The Gini coefficient cannot be inferred from Lorenz-style information.
\end{choiceslist}
\item Which of the following tax policies is most likely to reduce after-tax income inequality?
\begin{choiceslist}
\item A regressive sales tax with no transfers
\item A proportional income tax with no transfers
\item A progressive income tax combined with transfer payments to low-income households
\item A lump-sum tax on all households of the same dollar amount
\item An excise tax on necessities purchased mostly by low-income households
\end{choiceslist}
\item A general sales tax is typically described as regressive because
\begin{choiceslist}
\item high-income households spend a larger fraction of income on taxed goods
\item low-income households usually spend a larger fraction of income on consumption
\item the tax rate rises as income rises
\item the tax is levied only on luxury goods
\item the tax reduces consumer surplus for all households equally
\end{choiceslist}
\item Which of the following policies is most likely to reduce long-run income inequality by raising workers' human capital?
\begin{choiceslist}
\item Eliminating all taxes on inherited wealth
\item Reducing access to primary and secondary education
\item Subsidizing broad access to education and job training
\item Replacing progressive taxation with a flat head tax
\item Imposing a price floor in the market for luxury goods
\end{choiceslist}
\end{enumerate}
\end{multicols*}
\label{LastPage}
\end{document}
